\documentclass[11pt,oneside]{article}

\pdfminorversion=7
\usepackage{QuizStyle}

\hypersetup{
    pdftitle={2026Fall Quiz for MATH261 Discrete Mathematics},
    pdfauthor={Jungwoo Kim},
    pdfsubject={Quiz for MATH261 Discrete Mathematics},
    pdfcreator={LaTeX}
}

\begin{document}

\quiztitle{2026Fall Quiz}

\begin{problem}{1.1}
Let $X=\R$ and let $R=\varnothing$ be the void relation on $X$. Determine whether $R$ is reflexive, symmetric, and transitive, and hence whether $R$ is an equivalence relation. Justify each answer.
\end{problem}

\begin{solution}
The relation $R$ is not reflexive, since $(x,x)\notin R$ for every $x\in X$.

The relation $R$ is symmetric: the implication $xRy\Rightarrow yRx$ is vacuously true because $xRy$ never occurs. Likewise, $R$ is transitive because the premise $xRy\land yRz$ never occurs, so
\[
    (xRy\land yRz)\Rightarrow xRz
\]
is vacuously true.

Thus $R$ is symmetric and transitive but not reflexive. Therefore $R$ is not an equivalence relation.
\end{solution}

\begin{problem}{1.2}
For $n\in\Nzero$, prove that
\[
    \sum_{i=0}^{n}\binom{n}{i}^2
    =\binom{2n}{n}.
\]
\end{problem}

\begin{solution}[1]
Let $U$ and $V$ be disjoint sets with $|U|=|V|=n$. Count the $n$-element subsets of $U\cup V$ according to the number $i$ of selected elements from $U$. If exactly $i$ elements are chosen from $U$, then $n-i$ elements must be chosen from $V$, giving
\[
    \binom{n}{i}\binom{n}{n-i}
    =\binom{n}{i}^2
\]
choices. Summing over $i=0,\ldots,n$ counts all $n$-element subsets of the $2n$-element set $U\cup V$. Hence
\[
    \sum_{i=0}^{n}\binom{n}{i}^2
    =\binom{2n}{n}.
\]
\end{solution}

\begin{solution}[2]
By symmetry of binomial coefficients and Vandermonde's identity,
\[
    \sum_{i=0}^{n}\binom{n}{i}^2
    =\sum_{i=0}^{n}\binom{n}{i}\binom{n}{n-i}
    =\binom{2n}{n}.
\]
\end{solution}

\end{document}
